\documentclass{article}
\usepackage{amsmath,amssymb,mathrsfs}
\usepackage[margin=1in]{geometry}
\allowdisplaybreaks
\begin{document}

\section*{Affine dynamic-bid-price LP: normalized P1 and paired D1}

\begin{equation}
(\mathbf{P1})\qquad
z_{\mathrm{P1}}=\max_X
\sum_{t,\,\vec r\in\mathscr R_t,\,\vec u\in\mathscr U_{\vec r}}
\left(\sum_j p_{t,j}f_j u_j\right)X_{t,\vec r,\vec u}.
\tag{18}
\end{equation}
Subject to
\begin{equation}
\begin{split}
&\sum_{\vec r\in\mathscr R_t,\,\vec u\in\mathscr U_{\vec r}}
 r_iX_{t,\vec r,\vec u}\\
&\quad=
\begin{cases}
c_i, &\text{if }t=1,\\[2pt]
\displaystyle
\sum_{\vec r\in\mathscr R_{t-1},\,\vec u\in\mathscr U_{\vec r}}
\left(r_i-\sum_j p_{t-1,j}a_{i,j}u_j\right)
X_{t-1,\vec r,\vec u},
&\forall t=2,\ldots,\tau,
\end{cases}
\qquad\forall i,t,
\end{split}
\tag{19}
\end{equation}
\begin{equation}
\sum_{\vec r\in\mathscr R_t,\,\vec u\in\mathscr U_{\vec r}}
X_{t,\vec r,\vec u}=1
\qquad\forall t,
\tag{21}
\end{equation}
\[
X\geqslant0.
\]
Here the nonnegativity applies to every decision variable
\(X_{t,\vec r,\vec u}\), for every \(t\),
\(\vec r\in\mathscr R_t\), and \(\vec u\in\mathscr U_{\vec r}\).
Equation (21) is the author's stated replacement for (20); the superseded
flow-normalization equations are not imposed. The \(t=1\) part of (19)
is retained, although the author notes that it can be eliminated.

The objective maximizes approximate total expected revenue over the finite
horizon. Equation (19) balances the expected mass of each resource type.
With (21), the variables \(X_{t,\vec r,\vec u}\) can be interpreted as
state--action probabilities. As the author cautions, an \(X\) solution
does not in general yield an implementable randomized control policy:
its support need not specify an action for every attainable state.

\section*{Fixed data, indices, states, and admissible actions}

There are \(m\) resources and \(n\) product or customer classes. Throughout,
\[
i\in\{1,\ldots,m\},\qquad
j\in\{1,\ldots,n\},\qquad
t\in\{1,\ldots,\tau\}.
\]
The last period in which sales can occur is \(\tau\). Unqualified sums
and quantifiers over \(i,j,t\) use these ranges.

The fixed resource vector and consumption matrix are
\[
\vec c\equiv(c_i),\qquad A\equiv(a_{i,j}),
\]
where \(A\) is an \(m\times n\) matrix. For every \(i\), \(c_i>0\)
is the integer amount of resource \(i\) available at the beginning of
period 1. For every \(i,j\), \(a_{i,j}\) is the integer amount of
resource \(i\) required by a class \(j\) customer. The \(j\)th column
of \(A\) is denoted by \(A^j\). An accepted class \(j\) customer earns
the fixed revenue
\[
f_j>0\qquad\forall j.
\]
The selected formulation uses \(f_j\), not the optional time-dependent
fare variant mentioned in the paper.

At most one customer arrives in each period. A class \(j\) customer
arrives in period \(t\) with probability \(p_{t,j}\); no customer arrives
with probability
\[
1-\sum_j p_{t,j}.
\]
These are fixed arrival probabilities, not decision variables. No
strict-positivity assumption on all arrival probabilities is imposed
in the formulation of Section 4.1.

The resource-state set and its period-specific versions are
\[
\mathscr R\equiv
\left\{\vec r\in\mathbb Z_+^m:
 r_i\in\{0,1,\ldots,c_i\}\ \forall i\right\},
\qquad \vec r\in\mathscr R,
\]
\[
\mathscr R_t=
\begin{cases}
\{\vec c\},&\text{if }t=1,\\
\mathscr R,&\text{if }t=2,\ldots,\tau.
\end{cases}
\]
In period 1 the state is \(\vec r=\vec c\). These sets are not reduced
to the states reachable under a particular initial state and arrival
process.

An action \(\vec u\) specifies whether each class would be accepted
if it arrived: \(u_j=1\) means acceptance and \(u_j=0\) means rejection.
The admissible-action sets are
\[
\vec u\in\mathscr U_{\vec r}
=\left\{\vec u\in\{0,1\}^n:
\vec r\geqslant A^j u_j\ \forall j\right\}
\qquad\forall\vec r\in\mathscr R.
\]
Thus a request must be rejected if insufficient resources remain, but
may also be rejected when resources suffice. The vectors \(\vec r\)
and \(\vec u\) index the columns of P1; they are not additional
continuous decision vectors in that LP.

Given \(\vec r\) and \(\vec u\) in period \(t\), the one-step outcomes are
\[
\begin{array}{c|c|c}
\text{event probability}&\text{current reward}&\text{next state}\\ \hline
p_{t,j}\quad\text{for each }j&f_j u_j&\vec r-A^j u_j\\
1-\sum_jp_{t,j}&0&\vec r
\end{array}
\]
The random next-period resource vector is denoted by \(\vec R^{t+1}\),
and its resource-\(i\) component by \(R_i^{t+1}\). Conditional
expectations below use only these one-step probabilities.

\section*{Affine approximation and paired D1}

The value \(v_t(\vec r)\) is the expected cost-to-go over periods
\(t,\ldots,\tau\) starting from state \(\vec r\). The terminal value is
\[
v_{\tau+1}(\vec r)\equiv0\qquad\forall\vec r.
\]
The selected affine functional approximation is
\begin{equation}
v_t(\vec r)\approx\theta_t+\sum_i V_{t,i}r_i
\qquad\forall t,\vec r.
\tag{14}
\end{equation}
Here \(V_{t,i}\) estimates the marginal value of a unit of resource
\(i\) in period \(t\), and \(\theta_t\) is a constant offset with
respect to the resource state. They are the decision coefficients in
D1, not fixed input prices. The basis functions specializing the general
functional approximation are
\[
\phi_i(\vec r)=r_i\quad\forall i,
\qquad \phi_0(\vec r)=1.
\]
The terminal coefficients are fixed as stated in Section 4.1:
\[
\theta_{\tau+1}=0,\qquad V_{\tau+1,i}=0\quad\forall i.
\]

The source-identified dual counterpart is
\begin{equation}
(\mathbf{D1})\qquad
\min_{\theta,V}\ \theta_1+\sum_i V_{1,i}c_i,
\tag{16}
\end{equation}
subject to
\begin{equation}
\begin{split}
\theta_t-\theta_{t+1}
+\sum_i\left(V_{t,i}r_i
 -V_{t+1,i}\mathrm E[R_i^{t+1}\mid\vec r,\vec u]\right)
\geqslant\sum_j p_{t,j}f_j u_j\\
\forall t,\quad\vec r\in\mathscr R_t,\quad
\vec u\in\mathscr U_{\vec r}.
\end{split}
\tag{17}
\end{equation}
The coefficients \(\theta_t,V_{t,i}\) for \(t=1,\ldots,\tau\)
are unrestricted real LP variables; Section 4.1 imposes no sign or
monotonicity constraints on them.

The unnumbered expansion of (17) printed immediately after it is
\[
\begin{split}
\theta_t-\theta_{t+1}
+\sum_i\left(
V_{t,i}r_i-V_{t+1,i}
\left(r_i-\sum_j p_{t,j}a_{i,j}u_j\right)
\right)
\geqslant\sum_j p_{t,j}f_j u_j\\
\forall t,\quad\vec r\in\mathscr R_t,\quad
\vec u\in\mathscr U_{\vec r}.
\end{split}
\]
This is the same constraint, not an additional restriction. It also
specifies explicitly the conditional-expectation coefficient used
in (17).

\paragraph{Scope and source.}
The operative formulation is Section 4.1 of Daniel Adelman,
``Dynamic Bid Prices in Revenue Management,'' \emph{Operations Research}
55(4), 647--661 (2007), printed p.~652 (physical PDF p.~6), with the
problem data and state--action definitions on printed pp.~648--649
(physical PDF pp.~2--3). The general basis-function programs of
Section 3 specialize to the rows above; they introduce no additional
constraints for this affine version. The later structural results and
computational monotonicity restrictions are not imposed in this
selected Section 4.1 formulation. The appendix's quasistatic
approximation is a different formulation and is not included.

\end{document}
